\begin{abstract}
When new visual evidence contradicts an initial interpretation, do vision-language models (VLMs) genuinely revise their beliefs, or do they merely rationalize their first guess?  We introduce \emph{Evidence Update Prompting} (\method), a two-phase evaluation protocol inspired by defeasible and non-monotonic reasoning from cognitive science.  In \phaseA, a model receives limited pre-event evidence and forms an initial hypothesis; in \phaseB, additional post-event evidence arrives that often requires the model to revise.  We compare three prompting strategies---\emph{Baseline} (standard answer), \emph{Belief-State} (explicit hypothesis tracking with confidence), and \emph{Counterfactual Update} (``would your answer differ without the new evidence?'')---across three frontier VLMs (GPT-4o, Gemini~2.0~Flash, Claude~3.5~Sonnet) on 52 BlackSwan-style scenarios requiring abductive reasoning about surprising events.  Our findings reveal that (i)~all models exhibit substantial \emph{stubbornness}: 37--62\% of initially incorrect answers are never revised despite conflicting evidence; (ii)~Belief-State prompting reduces stubbornness by 13--18 percentage points and increases accuracy by 4--8 pp over baseline; (iii)~Counterfactual prompting helps models \emph{recognize} when evidence matters (59--63\% say ``yes, my answer would differ'') but produces only modest behavioral change; and (iv)~models display striking confidence inflation in \phaseB, with high-confidence predictions rising 2--3$\times$ regardless of whether the answer actually changed.  These results establish that current VLMs lack genuine belief revision mechanisms and instead engage in post-hoc rationalization, pointing toward architectures with explicit epistemic state tracking as a path forward.
\end{abstract}
